\documentclass[11pt]{amsart}
\usepackage{mathrsfs}


%standard packages
%------------------------------
%The following three packages are all included in amsart
\usepackage{amsmath}
\usepackage{amssymb} %to produce blackboard bold characters(\supsetneq)
%\usepackage{amsthm} %to define theorem-like structures
%-------------------------------
\usepackage{enumitem} %lay­out of the three ba­sic list en­vi­ron­ments
\usepackage{mathtools} %an extension package to amsmath + loading amsmath
\usepackage[table]{xcolor} %color of tables
\usepackage[all]{xy} %xymatrix commutative diagram
%\usepackage{tikz} %pictures
%\usepackage{tikz-cd}%implies
%\usepackage{geometry} %define page dimensions
\usepackage{indentfirst} %Indent first paragraph
%\usepackage{babel} %lan­guages
\usepackage{setspace} %set­ting the spac­ing be­tween lines

%setting of hyperlinks
\usepackage[colorlinks,linkcolor=red,anchorcolor=green,citecolor=blue]{hyperref} %color of hyperlinks
\hypersetup{linktocpage = true} %color of hyperlinks of TOC

%special packages and notations (rarely used)
\usepackage{rotating} %Rotate an image, table or paragraph
\def\no{n{\raisebox{0.2ex}{\textsuperscript{o}}}} %numbering


%various tables and their settings
%\usepackage{ytableau} %young tableau
\usepackage{longtable} %long table
\newcolumntype{M}[1]{>{\centering\arraybackslash}m{#1}} %define dimension for long stable

\setlength{\textheight}{8in}
\setlength{\textwidth}{5.87in} \setlength{\oddsidemargin}{0.3in}
\setlength{\evensidemargin}{0.3in}\setlength{\footskip}{0.3in}
\setlength{\headsep}{0.3in}


%setting of fonts
\usepackage{mathpazo}
\usepackage{mathrsfs}
\DeclareFontFamily{OMS}{rsfs}{\skewchar\font'60}
\DeclareFontShape{OMS}{rsfs}{m}{n}{<-5>rsfs5 <5-7>rsfs7 <7->rsfs10 }{}
\DeclareSymbolFont{rsfs}{OMS}{rsfs}{m}{n}
\DeclareSymbolFontAlphabet{\scr}{rsfs}
\DeclareSymbolFontAlphabet{\scr}{rsfs}

% a beautiful font (with XeLatex)
\usepackage[T1]{fontenc}
%\usepackage[utf8]{inputenc}
%\usepackage[no-math]{fontspec}
%\setromanfont[Ligatures={Common}, Numbers={OldStyle}, Variant=01]{Linux Libertine O}

%pagestyle
\pagestyle{plain}%header empty; the footer contains the page number.
%\geometry{left=1.5cm,right=1.5cm,top=1.5cm,bottom=1cm} %dimension of pages
\renewcommand{\baselinestretch}{1} %line spacing
%\setlength{\parindent}{0in} %indent
\sloppy %avoiding overfulls


%mathcal letters
\newcommand\cA{{\mathcal A}}
\newcommand\cB{{\mathcal B}}
\newcommand\cC{{\mathcal C}}
\newcommand\cD{{\mathcal D}}
\newcommand\cE{{\mathcal E}}
\newcommand\cF{{\mathcal F}}
\newcommand\cG{{\mathcal G}}
\newcommand\cH{{\mathcal H}}
\newcommand\cI{{\mathcal I}}
\newcommand\cK{{\mathcal K}}
\newcommand\cL{{\mathcal L}}
\newcommand\cN{{\mathcal N}}
\newcommand\cM{{\mathcal M}}
\newcommand\cO{{\mathcal O}}
\newcommand\cP{{\mathcal P}}
\newcommand\cQ{{\mathcal Q}}
\newcommand\cR{{\mathcal R}}
\newcommand\cS{{\mathcal S}}
\newcommand\cT{{\mathcal T}}
\newcommand\cU{{\mathcal U}}
\newcommand\cV{{\mathcal V}}
\newcommand\cW{{\mathcal W}}
\newcommand\cX{{\mathcal X}}

%mathbb letters
%\newcommand\bbB{{\mathbb B}}
%\newcommand\bbC{{\mathbb C}}
\newcommand\bbF{{\mathbb F}}
\newcommand\bbK{{\mathbb K}}
\newcommand\bbN{{\mathbb N}}
\newcommand\bbP{{\mathbb P}}
\newcommand\bbQ{{\mathbb Q}}
\newcommand\bbR{{\mathbb R}}
\newcommand\bbS{{\mathbb S}}
\newcommand\bbZ{{\mathbb Z}}

%mathscr letters
\newcommand\sA{{\mathscr A}}
\newcommand\sB{{\mathscr B}}
\newcommand\sD{{\mathscr D}}
\newcommand\sE{{\mathscr E}}
\newcommand\sF{{\mathscr F}}
\newcommand\sG{{\mathscr G}}
\newcommand\sI{{\mathscr I}}
\newcommand\sL{{\mathscr L}}
\newcommand\sM{{\mathscr M}}
\newcommand\sN{{\mathscr N}}
\newcommand\sO{{\mathscr O}}
\newcommand\sQ{{\mathscr Q}}
\newcommand\sT{{\mathscr T}}
\newcommand\sW{{\mathscr W}}
\newcommand\sH{{\mathscr H}}


%Special Abb.
\newcommand{\id}{{\rm id}}
\newcommand{\rk}{{\rm rk}}
\newcommand{\codim}{{\rm codim}}
\newcommand{\Cl}{{\rm Cl}}
\newcommand{\im}{{\rm im}}
\newcommand{\Bs}{{\rm Bs}}
\newcommand{\locus}{\rm locus}
\newcommand{\Sec}{\rm Sec}
\newcommand{\defect}{\rm def}
\newcommand{\Aut}{\rm Aut}
\newcommand{\Ext}{\rm Ext}
\newcommand{\Lag}{\rm Lag}
\newcommand{\Gr}{\rm{Gr}}
\newcommand{\Br}{{\rm Br}}



%Objects in math
\DeclareMathOperator*{\bs}{Bs}
\DeclareMathOperator*{\Irr}{Irr}
\DeclareMathOperator*{\Sym}{Sym}
\DeclareMathOperator*{\pic}{Pic}
\DeclareMathOperator*{\reg}{reg}
\DeclareMathOperator*{\Sing}{Sing}
\DeclareMathOperator*{\Pic0}{Pic^0}
\DeclareMathOperator*{\Nm}{Nm}
\DeclareMathOperator*{\red}{red}
\DeclareMathOperator*{\num}{num}
\DeclareMathOperator*{\norm}{norm}
\DeclareMathOperator*{\nons}{nons}
\DeclareMathOperator*{\lin}{lin}
\DeclareMathOperator*{\supp}{Supp}
\DeclareMathOperator*{\NS}{NS}
\DeclareMathOperator*{\Exc}{Exc}
\DeclareMathOperator*{\PGL}{PGL}

%Moduli spaces and objects attacted
\newcommand{\NEX}{\overline{\mbox{NE}}(X)}
\newcommand{\NAX}{\overline{\mbox{NA}}(X)}
\newcommand{\NE}[1]{\ensuremath{\overline{\mbox{NE}}(#1)}}
\newcommand{\NA}[1]{\ensuremath{\overline{\mbox{NA}}(#1)}}
\newcommand{\chow}[1]{\ensuremath{\mathcal{C}(#1)}}
\newcommand{\chown}[1]{\ensuremath{\mathcal{C}_n(#1)}}
\newcommand{\douady}[1]{\ensuremath{\mathcal{D}(#1)}}
\newcommand{\Chow}[1]{\ensuremath{\mbox{\rm Chow}(#1)}}
\newcommand{\Hilb}[1]{\ensuremath{\mbox{\rm Hilb}(#1)}}
\newcommand{\Univ}{\ensuremath{\mathcal{U}}}
\newcommand{\RC}[1]{\ensuremath{\mbox{\rm RatCurves}^n(#1)}}

%sheaves operators
\newcommand\Ctwo{\ensuremath{C^{(2)}}}
\newcommand{\Proj}[1]{\ensuremath{\mbox{\rm Proj}(#1)}}
\newcommand{\Hombir}[2]{\ensuremath{\mbox{$\rm{Hom}_1$}(#1,#2)}}
\newcommand{\cHom}[2]{\ensuremath{\mathcal{H}om_{\mathcal{O}_X}(#1,#2)}}
\newcommand{\sHom}[2]{\ensuremath{\mathscr{H}om_{\mathscr{O}_X}(#1,#2)}}
\newcommand{\cExt}{{\mathcal Ext}}
\newcommand{\cTor}{{\mathcal Tor}}
\newcommand{\defeq}{{\vcentcolon=}}

%math operators
\newcommand\bp{{\bar\partial}}

%Morphisms
\newcommand{\meom}[3]{\ensuremath{#1\colon #2\dashrightarrow #3}}
\newcommand{\morp}[3]{\ensuremath{#1\colon #2\rightarrow #3}}

%Theorem type newenviroment
% Theorem type environments
\theoremstyle{plain}
\newtheorem{thm}{Theorem}[section]
\newtheorem{lemma}[thm]{Lemma}
\newtheorem{prop}[thm]{Proposition}
\newtheorem{cor}[thm]{Corollary}
\newtheorem{defn}[thm]{Definition}
\newtheorem{claim}[thm]{Claim}
\newtheorem{conjecture}[thm]{Conjecture}
\newtheorem{question}[thm]{Question}
\newtheorem{problem}[thm]{Problem}

\theoremstyle{definition}
\newtheorem{example}[thm]{Example}
\newtheorem{remark}[thm]{Remark}
\newtheorem{notation}[thm]{Notation}
\newtheorem{convention}[thm]{Convention}




%%%%%%%%%%%%%%%%%%%%%%%%%



\newcommand{\btheorem}{\begin{thm}}
	\newcommand{\etheorem}{\end{thm}}
\newcommand{\bproposition}{\begin{prop}}
	\newcommand{\eproposition}{\end{prop}}
\newcommand{\bdefinition}{\begin{defn}}
	\newcommand{\edefinition}{\end{defn}}
\newcommand{\bcorollary}{\begin{cor}}
	\newcommand{\ecorollary}{\end{cor}}
\newcommand{\bproof}{\begin{proof}}
	\newcommand{\eproof}{\end{proof}}
\newcommand{\bremark}{\begin{remark}}
	\newcommand{\eremark}{\end{remark}}
\newcommand{\eexample}{\end{example}}
\newcommand{\bexample}{\begin{example}}

\newcommand{\elemma}{\end{lemma}}
\newcommand{\blemma}{\begin{lemma}}

\newcommand{\la}{\langle}
\newcommand{\ra}{\rangle}
\newcommand{\sq}{\sqrt{-1}}
\newcommand{\suml}{\sum\limits}
\newcommand{\p}{\partial}
\newcommand{\ord}{ord}
\newcommand{\Div}{ Div}

\renewcommand{\bar}{\overline}
\newcommand{\eps}{\varepsilon}
\newcommand{\pa}{\partial}
\renewcommand{\phi}{\varphi}
\newcommand{\OO}{\mathcal O}
\newcommand{\beq}{\begin{equation}}
\newcommand{\eeq}{\end{equation}}
\newcommand{\ee}{\end{eqnarray*}}
\newcommand{\be}{\begin{eqnarray*}}
\renewcommand{\red}{\textcolor[rgb]{1.00,0.00,0.00}}
\newcommand{\yel}{\textcolor[rgb]{0.00,0.00,1.00}}

\newcommand{\bd}{\begin{enumerate}}
	\newcommand{\ed}{\end{enumerate}}
\newcommand{\ul}{\underline}
\renewcommand{\hat}{\widehat}
\renewcommand{\tilde}{\widetilde}
\newcommand{\ch}{{ ch}}
\newcommand{\qtq}[1]{\quad\mbox{#1}\quad}
\renewcommand{\bp}{\bar{\partial}}
\newcommand{\Om}{\Omega}
\newcommand{\td}{ Td}
\newcommand{\ind}{Ind}
\newcommand{\ds}{\oplus}
\newcommand{\bds}{\bigoplus}
\newcommand{\ts}{\otimes}
\newcommand{\bts}{\bigotimes}
\newcommand{\diag}{ diag}
\newcommand{\st}{\stackrel}

\newcommand{\T}{{\mathbb T}}
%\newcommand{\U}{{\mathbb U}}
\newcommand{\V}{{\mathbb V}}
\newcommand{\W}{{\mathbb W}}
\renewcommand{\S}{{\mathbb S}}
\newcommand{\Y}{{\mathbb Y}}
\newcommand{\Z}{{\mathbb Z}}
\renewcommand{\>}{\rightarrow}


%mathfrak
\newcommand{\ssA}{{\mathfrak A}}
\newcommand{\ssB}{{\mathfrak B}}
\newcommand{\ssC}{{\mathfrak C}}
\newcommand{\ssD}{{\mathfrak D}}
\newcommand{\ssE}{{\mathfrak E}}
\newcommand{\ssF}{{\mathfrak F}}
\newcommand{\ssg}{{\mathfrak g}}
\newcommand{\ssH}{{\mathfrak H}}
\newcommand{\ssI}{{\mathfrak I}}
\newcommand{\ssJ}{{\mathfrak J}}
\newcommand{\ssK}{{\mathfrak K}}
\newcommand{\ssL}{{\mathfrak L}}
\newcommand{\ssM}{{\mathfrak M}}
\newcommand{\ssN}{{\mathfrak N}}
\newcommand{\ssO}{{\mathfrak O}}
\newcommand{\ssP}{{\mathfrak P}}
\newcommand{\ssQ}{{\mathfrak Q}}
\newcommand{\ssR}{{\mathfrak R}}
\newcommand{\ssS}{{\mathfrak S}}
\newcommand{\ssT}{{\mathfrak T}}
\newcommand{\ssU}{{\mathfrak U}}
\newcommand{\ssV}{{\mathfrak V}}
\newcommand{\ssW}{{\mathfrak W}}
\newcommand{\ssX}{{\mathfrak X}}
\newcommand{\ssY}{{\mathfrak Y}}
\newcommand{\ssZ}{{\mathfrak Z}}


\newcommand{\bbC}{\mathbb C}

\newcommand{\B}{{\mathbb B}}
\newcommand{\C}{{\mathbb C}}
\newcommand{\D}{{\mathbb D}}
\newcommand{\E}{{\mathbb E}}
\newcommand{\F}{{\mathbb F}}
%\newcommand{\G}{{\mathbb G}}
\renewcommand{\H}{{\mathbb H}}
%\newcommand{\I}{{\mathbb I}}
\newcommand{\J}{{\mathbb J}}
\newcommand{\K}{{\mathbb K}}
\renewcommand{\L}{{\mathbb L}}
\newcommand{\M}{{\mathbb M}}
\newcommand{\N}{{\mathbb N}}
\renewcommand{\P}{{\mathbb P}}
\newcommand{\Q}{{\mathbb Q}}
\newcommand{\R}{{\mathbb R}}
%\newcommand{\bS}{{\mathbb S}}


\newcommand{\ov}[1]{\overline{#1}}
\newcommand{\de}{\partial}
\newcommand{\db}{\overline{\partial}}
\newcommand{\mn}{\sqrt{-1}}
\newcommand{\ddbp}{\sqrt{-1}\partial \overline{\partial} \phi}
\newcommand{\ddbar}{\sqrt{-1} \partial \overline{\partial}}


%%%%%%%%%%%%%%%%%%%%%%%%%%%%%



%itemize&enumerate-->nonindent
%\setlist[itemize]{leftmargin=*}
%\setlist[enumerate]{leftmargin=*}

%numbering of lists
\renewcommand{\labelenumi}{(\arabic{enumi})} %numbering of enumerate
\renewcommand{\labelenumii}{(\theenumi.\arabic{enumii})} %subenumerate
\numberwithin{equation}{section} %numbering of equations
\renewcommand\thesubsection{\thesection.\Alph{subsection}} %subsection

%depth of TOC
\setcounter{tocdepth}{1} %include till subsections in TOC

%quotient
%\newcommand{\bigslant}[2]{{\raisebox{.2em}{$#1$}\left/\raisebox{-.2em}{$#2$}\right.}}

\makeatletter

\ifnum\@ptsize=0 \addtolength{\hoffset}{-0.3cm} \fi \ifnum\@ptsize=2 \addtolength{\hoffset}{0.5cm} \fi


%%%%%%%%%%%%%%%%%%%%%%%%%%%%%%%%%%%%%%%%%%%%%%%%%%%%%%%%%%%%%%%%%%%%%%%%%%%
\usepackage{fancyhdr}
\pagestyle{fancy}
% \pagestyle{myheadings} %\pagestyle{headings}
\renewcommand{\leftmark}{{\scriptsize Compact K\"ahler manifolds with quasi-positive second Chern-Ricci curvature}}
\renewcommand{\rightmark}{{\scriptsize
		Xiaokui Yang}}




%informations of paper
\title{Compact K\"ahler manifolds with quasi-positive second Chern-Ricci curvature}
%\date{\today}

\subjclass[2010]{53C55,14E08,32Q15}
\keywords{Chern-Ricci curvature, vanishing theorem, rational connectedness}




\author{Xiaokui Yang}

\address{Xiaokui Yang, Department of Mathematics and Yau Mathematical Sciences Center, Tsinghua University, Beijing, 100084, China}
\email{xkyang@mail.tsinghua.edu.cn}

\begin{document}
	
	\begin{abstract} Let $X$ be a compact K\"ahler manifold. We prove that if $X$ admits a smooth Hermitian metric $\omega$ with \emph{quasi-positive} second Chern-Ricci curvature $\mathrm{Ric}^{(2)}(\omega)$, then $X$ is projective and rationally connected. In particular, $X$ is simply connected.
	\end{abstract}
	
	
	\maketitle
	
	\tableofcontents
	
	\vspace{-0.2cm}
	
	\section{Introduction}
    The geometry of complex manifolds are characterized by various positivity notions in complex differential geometry and algebraic geometry. Since the seminal works of Mori   and Siu-Yau on the solutions to
	Hartshorne
	conjecture   and Frankel conjecture
	(\cite{Mori1979}, \cite{SiuYau1980}) on the characterizations of projective spaces,   many remarkable generalizations  have been
	established,   for instances, Mok's uniformization theorem on compact K\"ahler manifold
	with non-negative holomorphic bisectional curvature (\cite{Mok1988}) and  the  works of Campana,
	Demailly, Peternell and Schneider (\cite{CampanaPeternell1991},
	\cite{DemaillyPeternellSchneider1994})  on the
	structure of projective manifolds with nef tangent bundles. For more geometric characterizations, we refer to \cite{Mok1988,CampanaPeternell1991,DemaillyPeternellSchneider1994,DemaillyPeternellSchneider2001,Wu2002,Chau2012,LiWuZheng2013,Liu2014b,FengLiuWan2017,HuangLeeTamTong2018,NiZheng2019b,Liu2019,LiOuYang2019,LiuOuYang2020} and the references therein.\\
	
	    The  holomorphic sectional curvature also carries much geometric information of   complex manifolds. Indeed, thanks to the breakthrough work \cite{WuYau2016} of Wu-Yau, it is well-known that  a compact K\"ahler manifold $X$ with negative or quasi-negative holomorphic sectional curvature is algebraic and has ample canonical bundle (\cite{TosattiYang2017,WuYau2016a,Diverio2019}), which settles down a long-standing conjecture of S.-T. Yau affirmatively. For more recent works on  non-positive holomorphic sectional curvature, we refer to \cite{Wong1981,Heier2010,Heier2016,Heier2018,Nomura2018,Guenancia2018,YangZheng2019TAMS,WuYau2020} and the references therein. On the other hand, in his “Problem section”,
	    S.-T. Yau proposed the well-known conjecture \cite[Problem~47]{Yau1982} that compact K\"ahler manifolds with positive holomorphic sectional curvature must be projective and rationally connected. Recently, we solved this conjecture affirmatively in \cite{Yang2018} by introducing the concept of RC-positivity for abstract vector bundles, and many properties of such bundles are developed in \cite{Yang2018, Yang2018a,Yang2018b,Yang2021}. For instance, we proved in \cite{Yang2018b}  that a compact K\"ahler manifold with uniformly RC-positive tangent bundle must be projective and rationally connected.  By using the ideas of RC-positivity and some deep analytical techniques in algebraic geoemtry, Shin-ichi Matsumura established in \cite{Matsumura2018a,Matsumura2018b,Matsumura2018} a structure theorem for projective manifolds  with non-negative holomorphic sectional curvature, which is analogous to fundamental works in \cite{Mok1988, Campana2015a,Cao2017a,Cao2019} for manifolds with various non-negative properties (see also an approach in \cite{HeierWong2015}). In the same spirit, Lei Ni and Fangyang Zheng introduced in \cite{Ni2018,NiZheng2018a} various notions of Ricci curvature and scalar curvature to obtain rational connectedness of compact K\"ahler manifolds.\\
	
	
	
	In this paper, we investigate the geometry characterized by the first and second Chern-Ricci curvatures on Hermitian manifolds. Recall that for a Hermitian metric $\omega_g$,  its Chern curvature tensor has components $R_{i\bar j k\bar\ell}$. The first Chern-Ricci curvature is $$\mathrm{Ric}^{(1)}(\omega)=\sq \left(g^{k\bar\ell}R_{i\bar j k\bar \ell}\right)dz^i\wedge d\bar z^j=-\sq \p\bp\log\det(h_{k\bar\ell}) $$
	and the second Chern-Ricci curvature is 
	$$\mathrm{Ric}^{(2)}(\omega)=\sq \left(g^{k\bar\ell}R_{ k\bar \ell i\bar j}\right)dz^i\wedge d\bar z^j.$$
	When the Hermitian metric $\omega_g$ is not K\"ahler,  $\mathrm{Ric}^{(1)}(\omega)$ and $\mathrm{Ric}^{(2)}(\omega)$ are not necessarily the same. It is well-known that $\mathrm{Ric}^{(1)}(\omega)$ represents the first Chern class of the complex manifold. However, the geometry of $\mathrm{Ric}^{(2)}(\omega)$ is still mysterious.\\
	
	Thanks to the celebrated Calabi-Yau theorem (\cite{Yau1978}),  we know that  a
	compact K\"ahler manifold $X$ has a Hermitian metric with positive first Chern-Ricci curvature $\mathrm{Ric}^{(1)}(\omega)$
	if and only if $X$ is Fano. As an analog, we proved in \cite{Yang2018} that if a compact K\"ahler manifold
	admits a smooth Hermitian metric with positive second Chern-Ricci curvature $\mathrm{Ric}^{(2)}(\omega)$, then $X$ is projective and rationally connected. This result is also a generalization of the classical result of Campana \cite{Campana1992} and Koll\'ar-Miyaoka-Mori \cite{KollarMiyaokaMori1992} that Fano manifolds are rationally connected. The main result of this paper is the following theorem. 
	
	
	
	
	
		\btheorem\label{thm:quasipositivesecondchernriccirationallyconnected} Let $X$ be a compact K\"ahler manifold. If there exist a smooth Hermitian metric $\omega$ on $X$ and a smooth Hermitian metric $h$ on the holomorphic tangent bundle $TX$ such that  $$\mathrm{tr}_\omega R^{(TX,h)}\in \Gamma(X, \mathrm{End}(TX))$$ is quasi-positive. Then $X$ is projective and rationally connected. In particular, $X$ is simply connected.
	\etheorem
\noindent Here quasi-positive means non-negative everywhere and strictly positive at some point.  We follow the ideas in \cite{Campana2015a} and \cite{Graber2003} in the proof of Theorem \ref{thm:quasipositivesecondchernriccirationallyconnected}, and the key new ingredient is an integration argument  for singular Hermitian metrics instead of the pointwise maximum principle for RC-positive vector bundles employed in \cite{Yang2018}, since the latter does not work for manifolds with quasi-positive curvature tensors. It is easy to see that many compact K\"ahler manifolds with stable tangent bundle and positive slope can support smooth Hermitian metrics with positive or quasi-positive second Chern-Ricci curvature as required in Theorem \ref{thm:quasipositivesecondchernriccirationallyconnected} (e.g. \cite{UhlenbeckYau1986,Donaldson1987}). On the other hand, the K\"ahler condition in Theorem \ref{thm:quasipositivesecondchernriccirationallyconnected} is necessary (\cite[Section~6]{LiuYang2017}).
As a special case of Theorem \ref{thm:quasipositivesecondchernriccirationallyconnected}, one has



	\bcorollary \label{quasipositivesecondchernriccirationallyconnected} Let $X$ be a compact K\"ahler manifold. If $X$
	admits a smooth Hermitian metric $\omega$ with quasi-positive second Chern-Ricci curvature $\mathrm{Ric}^{(2)}(\omega)$, then $X$ is projective and rationally connected. In particular, $X$ is simply connected.
	\ecorollary
	


\noindent	By using  Corollary \ref{quasipositivesecondchernriccirationallyconnected}  and Yau's theorem \cite{Yau1978}, one has the following generalization of the result of  Campana \cite{Campana1992} and Koll\'ar-Miyaoka-Mori \cite{KollarMiyaokaMori1992}.
	
		\bcorollary \label{weakfano} Let $X$ be a compact K\"ahler manifold. If $X$
	admits a smooth Hermitian metric $\omega$ with quasi-positive first Chern-Ricci curvature $\mathrm{Ric}^{(1)}(\omega)$, then $X$ is projective and rationally connected. In particular, $X$ is simply connected.
	\ecorollary
	
	
	
	
	  \noindent Similarly,	for quasi-negative first Chern-Ricci curvature $\mathrm{Ric^{(1)}}(\omega)$, one has
	
	
	\btheorem \label{quasinegative} Let $(X,\omega)$ be a compact Hermitian manifold with quasi-negative first Chern-Ricci curvature $\mathrm{Ric^{(1)}}(\omega)$.
	If $X$ contains no rational curve, then $X$ is  projective  and $K_X$ is ample.
	\etheorem
	
		\noindent  This result is a straightforward consequence of deep results in complex analytical and algebraic geometry (\cite{Mori1982,Siu1984,Demailly1985,Kawamata1985,Birkar2010,Cascini2013}), which would be of independent interest from the viewpoint of complex differential geometry. It is known that  compact complex manifolds with quasi-negative (or quasi-positive) first Chern-Ricci curvature
	$\mathrm{Ric^{(1)}}(\omega)$ are Moishezon (\cite{Siu1984,Demailly1985}), which are not necessarily projective (e.g. \cite{Ma2007}).
	There are also many compact complex manifolds containing no rational curves, for instances,
	hyperbolic manifolds and Hermitian manifolds with non-positive holomorphic sectional curvature. Theorem \ref{quasinegative} also provides another proof of \cite[Corollary~1.1]{Lee2018} that  compact Hermitian manifolds with non-positive holomorphic bisectional curvature and quasi-negative first Chern-Ricci curvature are projective manifolds with ample canonical bundles,  which was established by purely analytical method.\\
	
	
	
	\noindent 
	The following conjecture on quasi-positive holomorphic sectional curvature is well-known and still widely open (e.g. \cite[Conjecture~1.9]{Yang2018b}) and in the special case when $X$ is projective it was confirmed affirmatively by Matsumura  (\cite{Matsumura2018b}). 
	
	\begin{conjecture}\label{quasipositiveHSC} Let $(X,\omega)$ be a compact K\"ahler manifold. If it has quasi-positive holomorphic sectional curvature, then $X$ is projective and rationally connected.
	\end{conjecture}
	
	
	

	
	
	\noindent{\bf Acknowledgements.}  The author would like to thank Jie Liu, Wenhao Ou, Valentino Tosatti and S.-T. Yau for inspiring discussions and useful communications. 
	
	
	
	\vskip 2\baselineskip
	


	\section{Singular Hermitian metrics on vector bundles}
	
	Let $X$ be a complex manifold and $\omega$ be a smooth Hermitian metric on $X$. Locally, we can write the curvature tensor of $(T_X, \omega)$ as 
	$$R_{i\bar j k\bar \ell}=-\frac{\p^2 g_{k\bar\ell}}{\p z^i\p\bar z^j}+g^{p\bar q}\frac{\p g_{k\bar q}}{\p z^i}\frac{\p g_{p\bar \ell}}{\p\bar z^j}.$$ where $\omega=\sq g_{i\bar j}dz^i\wedge d\bar z^j$.
	The \emph{first Chern-Ricci curvature} $\mathrm{Ric}^{(1)}(\omega)=\sq R_{i\bar j}dz^i\wedge d\bar z^j$  has components
	$R^{(1)}_{i\bar j}=g^{k\bar\ell}R_{i\bar j k\bar \ell}=-\frac{\p^2\log \det (g_{k\bar\ell})}{\p z^i\p\bar z^j}.$
	The \emph{second Chern-Ricci curvature} is  $\mathrm{Ric}^{(2)}(\omega)=\sq R^{(2)}_{i\bar j}dz^i\wedge d\bar z^j$ where
	$R^{(2)}_{i\bar j}=g^{k\bar\ell}R_{ k\bar \ell i\bar j}$. The scalar curvature $s$ of the Chern connection is $\mathrm{tr}_\omega \mathrm{Ric^{(1)}}(\omega)$ which is also the same as $\mathrm{tr}_\omega \mathrm{Ric^{(2)}}(\omega)$. 
	
		Let $E\>X$ be a holomorphic vector bundle and $h$ be a smooth Hermitian metric on $E$. The curvature $R^E$ of the Chern connection $\nabla $ on $(E,h)$ has a similar formula $$R^E_{i\bar j \alpha\bar\beta}=-\frac{\p^2 h_{\alpha\bar\beta}}{\p z^i\p\bar z^j}+h^{\gamma\bar\delta}\frac{\p h_{\alpha\bar \delta}}{\p z^i}\frac{\p h_{\gamma\bar\beta}}{\p \bar z^j}.$$
	We define 
	$R^{(1)}_{i\bar j}=h^{\alpha\bar\beta} R_{i\bar j \alpha\bar\beta}$ and $R^{(2)}_{\alpha\bar\beta}=g^{i\bar j}R_{i\bar j \alpha\bar\beta}$.
	We also call $\mathrm{tr}_hR^E=\sq R_{i\bar j}^{(1)}dz^i\wedge d\bar z^j$ and $\mathrm{tr}_\omega R^E= \sq R_{\alpha\bar \beta}^{(2)}e^\alpha \ts \bar e^\beta$ the first and second Chern-Ricci curvature of $(E,h)$ with respect to the Hermitian manifold $(X,\omega)$ respectively. When $(E,h)=(TX,\omega)$, they are exactly the same as those curvatures of $(X,\omega)$.   A smooth Hermitian $(1,1)$-form $A=\sq A_{i\bar j}dz^i\wedge d\bar z^j$ on $X$ is call \emph{quasi-positive}, if $(A_{i\bar j})$ is non-negative everywhere and positive at some point of $X$. Similarly, we can define it for a  tensor $A \in \Gamma(X,\mathrm{End}(E))$.\\
	
	 Singular Hermitian metrics on line bundles are introduced in \cite{Demailly1992a} by Demailly. Let $L$ be a holomorphic vector bundle. A singular metric $h^L$ on $L$ can be written locally as $h^L=e^{-\phi}$ for some $\phi\in L^1_{\mathrm{Loc}}(X,\R)$, and the curvature $R^L=-\sq\p\bp\log h^L$ is defined in the sense of distributions. For singular Hermitian metrics on vector bundles, the definition would be very subtle, and we refer to \cite[Section~2]{Paun2018a} for a detailed discussion (see also \cite{Cataldo1998,Raufi2015,DengWangZhangZhou2020}). Recall that, for a smooth Hermitian metric $h$ on $E$, its curvature also takes the local form $R^E=\bp(h^{-1}\p h)$. For the specified purpose in this paper, we only consider singular metrics such that $h^{-1}\p h$ is locally integrable, and then we can use standard theory on distributions (e.g. \cite{Demailly2012}) to define the notion of weak positivity. 
	 
	 \bdefinition \label{definitionforsingularmetric} Let $X$ be a  complex manifold. A vector bundle $E$ is called to have positive second Chern-Ricci curvature in the sense of distributions, if there exist a smooth metric $\omega$ on $X$ and a singular Hermitian metric $h^E$ on $E$ such that
	 \beq \mathrm{tr}_\omega R^E \in \Gamma(X,\mathrm{End}(E))\eeq 
	 is strictly positive in the sense of distributions. The definition for non-negativity, quasi-positivity and etc. can be defined similarly.
	 \edefinition
	 
	 \bremark Of course, the notion of second Chern-Ricci curvature can be defined in a broader way by using similar constructions as in \cite{Cataldo1998}.  So far, it is not clear to the author whether the notions of Griffiths positivity or Nakano positivity for singular metrics defined in  \cite{Paun2018a,Raufi2015,DengWangZhangZhou2020} can imply the positivity of the second Chern-Ricci curvature in suitable sense, though it is obvious for smooth Hermitian metrics.
	\eremark
	
	\vskip 2\baselineskip
	
	\section{Vanishing theorems for singular Hermitian metrics}
	
	The main result of this section is the following theorem.
	
	
	\btheorem\label{vanishingforquasipositive} Let $E$ be a holomorphic vector bundle over a compact complex manifold $X$.
	Suppose there exist a smooth Hermitian metric $\omega$ on $X$ and a smooth Hermitian metric $h$ on  $E$ such that  $$\mathrm{tr}_\omega R^{(E,h)}\in \Gamma(X, \mathrm{End}(TX))$$
	is quasi-positive. We have the following assertions.
	
	\bd \item Any invertible \textbf{subsheaf} $L$ of $\sO(\ts^k E^*)$ $(k\geq 1)$ is not pseudo-effective.
	\item $\det E^*$ is not pseudo-effective. \ed
	\etheorem
	
	
	\noindent The key difficulty in the proof of Theorem \ref{vanishingforquasipositive} is to deal with a line bundle $L$ which is only a subsheaf of $\sO(\ts^k E^*)$. If it is indeed a subbundle, then its follows from a simple observation (c.f. \cite{Campana2015a}).
	
	
	\blemma\label{monotonicityofsecondchernricci} The second Chern-Ricci curvature is decreasing in subbundles and increasing in quotient bundles.
	\elemma
	
	
	\noindent It is well-known that the first Chern-Ricci curvature is not necessarily monotone as described in Lemma \ref{monotonicityofsecondchernricci}.
	The proof follows from a standard computation  and we include details here since we need it for the distribution case.
	
	\bproof  Let $(E,h)$ be a Hermitian vector bundle and $S$ be a holomorphic subbundle of $E$.
	Let $r$ be the rank of $E$ and $s$ the rank of $S$. Without
	loss of generality, we can assume,  at a fixed point $p\in X$, there
	exists a local holomorphic frame $\{e_1,\cdots,e_r\}$ of $E$
	centered at point $p$ such that $\{e_1,\cdots, e_s\}$ is a local
	holomorphic frame of $S$. Moreover, we can assume that
	$h(e_\alpha,e_\beta)(p)=\delta_{\alpha\beta}, \text{for } 1\leq
	\alpha, \beta \leq r.$ Hence, the curvature tensor of
	$S$ at point
	$p$ is $$ R^S_{i\bar j \alpha\bar\beta}=-\frac{\p^2
		h_{\alpha\bar\beta}}{\p z^i\p\bar z^j}+\sum_{\gamma=1}^{s}\frac{\p
		h_{\alpha\bar\gamma}}{\p z^i}\frac{\p h_{\gamma\bar\beta}}{\p\bar
		z^j}$$ where $1\leq \alpha,\beta\leq s$. For any Hermitian metric $\omega=\sq g_{i\bar j}dz^i\wedge d\bar z^j$ on $X$, we have
	\beq (\mathrm{tr}_\omega R^E)|_{S}-\mathrm{tr}_\omega R^S= \mathrm{tr}_\omega(R^E|_{S})-\mathrm{tr}_\omega R^S=\sq
	\sum_{\alpha,\beta=1}^sg^{i\bar j}\left(\sum_{\gamma=s+1}^r\frac{\p
		h_{\alpha\bar\gamma}}{\p z^i}\frac{\p h_{\gamma\bar\beta}}{\p\bar
		z^j}\right) e^\alpha\ts
	\bar e^\beta. \label{secondchernriccimonotone}\eeq
	It is easy to see that the right hand side of (\ref{secondchernriccimonotone}) is non-negative. The proof for quotient bundles is similar.
	\eproof
	
	\bcorollary\label{functorialforsecondricci} Let $E$ be a holomorphic vector bundle over a complex manifold $X$.
	
	\bd \item  If $E$ has  positive (resp, non-negative) second Chern-Ricci curvature, then each quotient bundle has    positive (resp, non-negative) second Chern-Ricci curvature.
	
	\item If $E$ has negative (resp, non-positive) second Chern-Ricci curvature, then each subbundle has   negative (resp, non-positive) second Chern-Ricci curvature.
	
	\item If $E$ has  positive (resp, non-negative, quasi-positive,  negative, non-positive, quasi-negative) second Chern-Ricci curvature, then so is $\ts^k E$ for each $k\geq 1$.
	
	\item  If $E$ has  positive (resp, non-negative, quasi-positive,  negative, non-positive, quasi-negative) second Chern-Ricci curvature, then so is $\mathrm{Sym}^{\ts k} E$ $(k\geq 1)$ and $\Lambda^pE$ $ (1\leq p\leq \rk(E))$.
	
	
	\ed
	\ecorollary
	
	\bproof We only need to prove $(3)$. From the expression of the induced curvature formula of
	$(\ts ^kE^{},\ts ^kh^{})$, one has
	$$R^{(\ts ^kE^{},\ts ^kh^{})}=\ts ^k R^{(E,h)}\in \Gamma\left(X,\Lambda^{1,1}T^*X\ts \mathrm{End}(\ts^kE)\right),$$
	and
	$$\mathrm{tr}_\omega R^{(\ts ^kE^{},\ts ^kh^{})}=\mathrm{tr}_\omega\left(\ts ^k R^{(E,h)}\right)=\ts^k\left(\mathrm{tr}_\omega R^{(E,h)}\right)\in \Gamma\left(X,\mathrm{End}(\ts^k E)\right).$$
	Hence, the result follows.
	\eproof
	
	
	\bremark By standard distribution theory, one has similar results as in   Lemma \ref{monotonicityofsecondchernricci}
and Corollary \ref{functorialforsecondricci}  for singular Hermitian metrics  in  Definition \ref{definitionforsingularmetric}. 	
	\eremark
	
	
	\blemma\label{nonwherevanishing} Suppose that $E$ has non-negative second Chern-Ricci curvature in the  sense of distributions and $L$ is a pseudo-effective line bundle.
	If $\sigma\in H^0(X,E^*\ts L^*)$ is a non-trivial holomorphic section, then $\sigma$ does not vanish everywhere.
	\elemma
	
	\bproof Let $\omega$ be a smooth Hermitian metric on $X$ and $h^E$ be a singular metric such that $\mathrm{tr}_\omega R^E$ is non-negative in the sense of distributions.
	Since $$\mathrm{tr}_{e^f\omega} R^{E}=e^{-f}\mathrm{tr}_\omega R^E,$$
	we can assume further that $\omega$ is a Gauduchon metric (\cite{Gauduchon1977}), i.e. $\p\bp\omega^{n-1}=0$
	where $\dim X=n$. Let $h^L$ be a singular metric on $L$ such that its curvature $\theta=-\sq\p\bp\log h^L\geq 0$
	in the sense of distributions. For any $\sigma\in H^0(X,E^*\ts L^*)$, we have
	\beq\mathrm{tr}_\omega\left(\sq \p\bp |\sigma|^2_{h^{E^*}\ts h^{L^*}}\right)
	=|\nabla \sigma|^2_{\omega,h^{E^*},h^{L^*}}+(\mathrm{tr}_\omega R^E\cdot h^{L^*}+\mathrm{tr}_\omega\theta\cdot h^{E^*})(\sigma,\sigma).\label{bochnerformula}\eeq
	Note that we are dealing with positive currents, and the product makes sense.
	An integration by part argument with respect to the Gauduchon metric $\omega$  shows that $\nabla\sigma=0$ a.e., and
	\beq (\mathrm{tr}_\omega\sq\p\bp) |\sigma|^2_{h^{E^*}\ts h^{L^*}}\geq 0 \eeq
	in the sense of distributions.  We know $|\sigma|^2_{h^{E^*}\ts h^{L^*}}$ is a global constant on the compact base $X$.
	If $\sigma$ is non-trivial, this constant is nonzero, i.e. $\sigma$ does not vanish everywhere.
	\eproof
	
	\bcorollary\label{subsheafsubbundle} Suppose that $E$ has non-negative second Chern-Ricci curvature  in the  sense of distributions, and  $L$ is an invertible
	subsheaf of $\sO(E^*)$. If $L$ is pseudo-effective, then $L$ is a line subbundle of $E^*$.
	\ecorollary
	
	\bproof Note that the subsheaf morphism $f: L\> \sO(E^*)$ induces a nonzero section $\sigma \in H^0(X,E^*\ts L^*)$. By Lemma \ref{nonwherevanishing},
	$\sigma$ does not vanish everywhere. This gives a trivial line subbundle  $\underline{\C}$ of the vector bundle $E^*\ts L^*$, and so $L$ is a line subbundle of $E^*$.
	\eproof
	
	
	
	\btheorem\label{distributionvanishing} Suppose that $E$ has quasi-positive second Chern-Ricci curvature  in the  sense of distributions, and  $L$ is a pseudo-effective line bundle. Then
	\beq H^0(X,E^*\ts L^*)=0.\eeq
	\etheorem
	
	\bproof By using a similar argument as in the proof of Lemma \ref{nonwherevanishing}, for any holomorphic section $\sigma\in H^0(X,E^*\ts L^*)$,
	we deduce  from (\ref{bochnerformula}) that the holomorphic section $\sigma$ vanishes on an open subset of $X$. By Aronszajn's principle (\cite{Aronszajn1957}), $\sigma$ is a zero section.
	\eproof
	
	
	

	
	
\begin{proof}[Proof of Theorem \ref{vanishingforquasipositive}]
	Let $L$ be an invertible subsheaf of $\sO(\ts^k E^*)$ for some
	$k\geq 1$. Suppose-to the contrary-that $L$ is pseudo-effective. By
	Corollary \ref{functorialforsecondricci}, $\ts^k E$ has
	quasi-positive second Chern-Ricci curvature. By Corollary
	\ref{subsheafsubbundle}, the pseudo-effective invertible subsheaf
	$L$ is actually a line subbundle of $\ts^k E^*$. It is easy to see
	that if $\ts^k E$ has quasi-positive second Chern-Ricci curvature,
	then its dual bundle  $\ts^k E^*$ has quasi-negative second
	Chern-Ricci curvature. By decreasing principle in Lemma
	\ref{monotonicityofsecondchernricci}, the second Chern-Ricci
	curvature of $L$ is quasi-negative. Since $L$ is a line bundle,
	that means, the induced metric $h^L$  has quasi-negative Chern
	scalar curvature $s=\mathrm{tr}_\omega (-\sq\p\bp\log h^L)$. Without
	loss of generality, we can assume $\omega$ is Gauduchon. Therefore \beq
	\int_X s\omega^n=n \int_X c_1^{\mathrm{BC}}(L)\wedge
	\omega^{n-1}<0.\eeq  Hence, $L$ can not be
	pseudo-effective (\cite[Proposition~3.1]{Yang2019CM}, ). This is a contradiction.
	
	On the other hand, since $\det E^*$ is a subbundle of $\ts^k E^*$ for some large $k$, we deduce $\det E^*$ can not be pseudo-effective. \end{proof}
	
	
	
	\vskip 2\baselineskip
	
	
	\section{The proof of main theorems}
	
	
	
In this section, we prove Theorem \ref{thm:quasipositivesecondchernriccirationallyconnected}, Corollary \ref{weakfano} and Theorem \ref{quasinegative}.
	
	\bproof[{Proof of Theorem \ref{thm:quasipositivesecondchernriccirationallyconnected}}]  We set $(E,h)=(TX,h)$.  By Corollary \ref{functorialforsecondricci} and Theorem \ref{distributionvanishing}, we have
	$$H^0(X,(T^*X)^{\ts k})=0.$$
	It is well-known that for any $p$ satisfying $1\leq p\leq \dim X$, $\Lambda^pT^*X$ is a direct summand of $(T^*X)^{\ts k}$ for some large $k$. Therefore, we have
	$$H_{\bp}^{p,0}(X,\C)\cong H^{0}(X,\Lambda^p T^*X)=0.$$
	In particular,  $H_{\bp}^{2,0}(X,\C)=0$. Since $X$ is K\"ahler, we deduce that $X$ is projective (\cite{Kodaira1954}).
	Now we following the ideas in \cite{Graber2003,Peternell2006,Campana2015a} to get the rational connectedness.
	By using Theorem \ref{vanishingforquasipositive}, we conclude that $K_X=\det T^*X$ is not pseudo-effective.
	Thanks to \cite{BoucksomDemaillyPuaunPeternell2013}, $X$ is actually uniruled.  Let $\pi:X\dashrightarrow Z$ be the
	associated MRC fibration of $X$. After possibly resolving the
	singularities of $\pi$ and $Z$, we may assume that $\pi$ is a proper
	morphism and $Z$ is smooth. By \cite[Corollary~1.4]{Graber2003}, it
	follows that the target of the MRC fibration is either a point or a
	positive dimensional variety which is not uniruled. Suppose $X$ is
	not rationally connected, then $\dim Z\geq 1$ and $Z$ is not
	uniruled. By using \cite{BoucksomDemaillyPuaunPeternell2013} again, $K_Z$ is indeed pseudo-effective.
	Since $K_Z=\det(T^*Z)$ is  a direct summand of the vector bundle $(T^*Z)^{\ts k}$ for some large $k$ and $\sO((T^*Z)^{\ts k})$ is a subsheaf of $ \sO\left((T^*X)^{\ts k}\right)$,  we obtain a pseudo-effective invertible subsheaf $K_Z$ of $\sO((T^*X)^{\ts k})$, which contradicts to part $(1)$ of Theorem \ref{vanishingforquasipositive}.
	\eproof
	


\bproof[Proof of Corollary \ref{weakfano}] By the celebrated Calabi-Yau theorem, there exists a smooth \emph{K\"ahler metric} $\omega_0$ on $X$ such that 
\beq \mathrm{Ric}^{(1)}(\omega_0)=\mathrm{Ric}^{(1)}(\omega).\eeq
Since $\omega_0$ is K\"ahler, one can deduce
$$\mathrm{Ric}^{(2)}(\omega_0)=\mathrm{Ric}^{(1)}(\omega_0).$$
Therefore, Corollary \ref{weakfano} follows from Corollary \ref{quasipositivesecondchernriccirationallyconnected}.
\eproof
	
	\bproof[Proof of Theorem \ref{quasinegative}] By using the criterion of Siu-Demailly (\cite{Siu1984,Demailly1985}), we know $X$ is a Moishezon manifold and $K_X$ is  big. On the other hand, by the deep result  \cite[Corollary~1.4.6]{Birkar2010}, it is proved in \cite[Theorem~3.1]{Cascini2013} that a Moishezon manifold without rational curve must be projective. Moreover, by the base-point-free theorem and the relative cone theorem \cite{Mori1982,Kawamata1985,KollarMiyaokaMori1992}, one can deduce that $K_X$ is ample since $X$ is a projective manifold of general type  without rational curve.
	\eproof
	


\vskip 2\baselineskip

	
	\def\cprime{$'$} %newcommand of prime over letters
	


	\renewcommand\refname{Reference}
	\bibliographystyle{alpha}
	\bibliography{RC2}
	
\end{document}
